\section*{Resumen}

En el siguiente reporte se presenta el trabajo realizado durante la práctica profesional II bajo la supervisión de la Dra. Andrea Colins, cuyo objetivo es analizar datos de la actividad neuronal del ganglio pedal de \textit{Aplysia} en un conjunto de animales jóvenes y viejos, además de aplicar métodos estadísticos para comparar la actividad neuronal entre ambos grupos.\\

Durante el desarrollo de la práctica se logra comprender el ritmo del trabajo de investigación en la universidad, la importancia de la organización, honestidad en los resultados y buenas prácticas de programación \\

Como resultado, se logra un acercamiento al trabajo de investigación de ingeniero en biotecnología en una universidad y 

. Finalmente, luego de aprender el proceso, se detectan limitaciones y se proponen recomendaciones, tales como sugerir una forma distinta de organización de muestras para evitar su acumulación.



